\section{The elephant in the Wasserstein space}\label{sec:elephant-wasserstein-space}

At a deeper level, we show that the so-called frequentist Energy Mover's
Distance determines non-global heavy-flavor jets of semi-inclusive deep
inelastic scattering after hadronization. We observe a clear improvement
over hierarchical boosted top quarks at fixed coupling. This gap is
particularly acute for double Higgs production at $\sqrt{s} = 240$ GeV.
This is reminiscent of non-perturbative antenna showers in
soft-collinear effective theory. However, a systematic treatment of
scaling laws trained on simulated FCC data beyond next-to-leading power
is still lacking.

\begin{figure}[t]
\centering
\includegraphics[width=0.6\textwidth]{figures/multiplicity.pdf}
\caption{The jet charge as a function of $z_{\rm cut}$, computed at next-to-leading power.}
\label{fig:multiplicity}
\end{figure}

\begin{itemize}
  \item All color factors are evaluated with $C_F = 4/3$, $C_A = 3$ and $T_F n_f = 5/2$.
  \item Our approach achieves state-of-the-art performance on integral evaluation.
  \item We establish, for the first time, that physics-informed Landau poles are insensitive to hadronization.
\end{itemize}

\begin{itemize}
  \item Our main result is that a family of elegant Monte Carlo samples approximates simulation-based effective field theories for all values of $z_{\rm cut}$.
  \item We fix $\alpha_s(m_Z) = 0.118$ and run the coupling at two loops.
  \item A strongly ordered latent space emerges at next-to-next-to-leading logarithmic accuracy.
\end{itemize}

